\documentclass[11pt]{article}

\usepackage[margin=1in]{geometry}
\usepackage{times}
\usepackage[T1]{fontenc}

\title{LeanAide: Extensible Autoformalization by Combinator Code Generation}
\author{
  Siddhartha Gadgil \and
  Anirudh Gupta \and
  Ajay Kumar Nair \and
  Malhar Alpesh Patel \and
  Swarnava Chakraborty
}
\date{}

\begin{document}
\maketitle

\begin{center}
\emph{Department of Mathematics, Indian Institute of Science, Bangalore \\
(Malhar Alpesh Patel: Indian Institute of Science Education and Research, Pune)}
\end{center}

\begin{abstract}
Autoformalization is the translation of natural language mathematics into a
formal language that is machine-verifiable. We present LeanAide, an
extensible, modular framework for autoformalization using large language
models, metaprogramming and proof automation in Lean. The main contribution
of this work is the design and implementation of a novel autoformalization
framework based on \emph{Combinator Code Generation}, where we combine
handlers for different mathematical constructs in a manner analogous to
Combinator Parsing.

While there are powerful autoformalization systems for Lean, we believe that
the modularity, extensibility, deep integration with Lean (including for
interfaces) and open source nature of LeanAide gives it advantages over
existing systems in many contexts. In particular, advances in both proof
automation (such as the \texttt{grind} tactic) and large language models
automatically translate to improvements in LeanAide. Further, users can
easily add support for new mathematical constructs by writing new handlers
and combining them with existing ones.
\end{abstract}

\noindent\textbf{Funding.} We gratefully acknowledge funding from Renaissance
Philanthropy under their AI for Math Fund.

\end{document}
